''Путешествие во времени'' в RR позволяет разработчику установить состояние
программы в любой момент, зафиксированный во время сессии записи. Для этого RR
всегда начинает воспроизведение с самого начала записанной трассы и быстро
''проматывает'' исполнение, используя свой журнал событий и значения ''внутренних
часов'', до указанной пользователем точки в прошлом. Достигнув цели, программа
останавливается в состоянии, идентичном тому, что было при оригинальной записи,
предоставляя полный доступ для анализа в отладчике.

Концепция Запись и Воспроизведение, лежащая в основе RR, стала основой для
реализации таких возможностей, как \textbf{Time-travel debugging} и
\textbf{Reverse Debugging}.

% RR позволяет достичь детерминированности, записывая недетерминированные
% события, используя ''внутренние часы'' и воспроизводя по журналу.

Данные способы отладки являются натуральным продолжением RR. ''Шаг вперед'' —--
это просто воспроизведение до нужной точки, а ''шаг назад'' достигается не
истинным обратным исполнением, а восстановлением состояния программы из ранее
сохраненного \textbf{чекпоинта} (снимка памяти и регистров) и быстрым
воспроизведением вперед с этого чекпоинта до точки, непосредственно
предшествующей желаемому моменту отката.
